\setstretch{1.2} 
\title{Policy Learning for Optimal Dynamic Treatment Regimes with Observational Data} 
%\thanks{The full version of the paper is available at \url{https://arxiv.org/abs/2106.12886}.}


\author{Shosei Sakaguchi\thanks{Faculty of Economics, The University of Tokyo, 7-3-1 Hongo, Bunkyo-ku, Tokyo 113-0033, Japan. Email: sakaguchi@e.u-tokyo.ac.jp.} }
%\author{}
\date{\today}
%\date{\today \medskip \\ \textit{Preliminary Draft}}

\maketitle
\thispagestyle{empty}

\vspace{-1cm}

\begin{abstract}
\begin{spacing}{1.2}
   Public policies and medical interventions often involve dynamic treatment assignments, in which individuals receive a sequence of interventions over multiple stages. We study the statistical learning of optimal dynamic treatment regimes (DTRs) that determine the optimal treatment assignment for each individual at each stage based on their evolving history. We propose a novel, doubly robust, classification-based method for learning the optimal DTR from observational data under the sequential ignorability assumption. The method proceeds via backward induction: at each stage, it constructs and maximizes an augmented inverse probability weighting (AIPW) estimator of the policy value function to learn the optimal stage-specific policy. We show that the resulting DTR achieves an optimal convergence rate of $n^{-1/2}$ for welfare regret under mild convergence conditions on estimators of the nuisance components.
    \bigskip \\
\noindent 
\textbf{Keywords:} Backward induction; Double robustness; Dynamic treatment; Policy learning; Sequential ignorability. \\
%\textbf{JEL codes:} C22, C44, C54
\end{spacing}
\end{abstract}

%%%%%%%%%%%%%%%%%%%%%%%%%%%%%%%%%%%%%%%%%%%%%%%%
%%%%%%%%%%%%%%%%%%%%%%%%%%%%%%%%%%%%%%%%%%%%%%%%
\thispagestyle{empty}
\clearpage
%\addtocounter{page}{-1}
%\newpage
\setstretch{1.5} 


\section{Introduction} \label{sec:introduction}


Public policies and medical interventions often involve dynamics in their treatment assignments. For example, some job training programs offer training sessions over multiple stages \citep[e.g.,][]{Lechner_2009,Rodriguez_et_al_2022}. In clinical medicine, physicians sequentially administer treatments in response to patients’ evolving conditions \citep[e.g.,][]{Wang_et_al_2012,Pelham_et_al_2016}. Multi-stage treatment assignments are also common educational programs spanning multiple grades \citep[e.g.,][]{krueger_1999,Ding_Lehrer_2010} and in dynamic marketing strategies \citep[e.g.,][]{Liu_2023}.

This study focuses on the optimal allocation of sequential treatments \citep{Robins_1986}, in which individuals receive interventions over multiple stages. In this context, treatment effects at each stage are often heterogeneous depending on prior treatments and related states. Thus, adapting sequential treatment allocation to evolving information can substantially enhance the welfare gains of multi-stage interventions.

We study statistical learning for optimal sequential treatment assignment using data from quasi-experimental or observational studies. Throughout this paper, we assume sequential ignorability \citep{Robins_1997}, meaning that treatment assignment at each stage is independent of potential outcomes, conditional on the history of prior treatments and observed states. Under this assumption, we propose a novel method for learning the optimal dynamic treatment regime (DTR), a sequence of stage-specific policies that determines the optimal treatment for each individual at each stage based on their history up to that point.


In developing our approach, we build on recent advances in doubly robust policy learning \citep{Athey_Wager_2020, zhou2022offline} and extend them to dynamic settings. We propose a doubly robust, classification-based method for learning the optimal DTR via backward induction, which sequentially estimates the optimal policy from the final stage to the first. At each step, the method constructs an augmented inverse probability weighting (AIPW) estimator of the policy value function by combining estimators of the propensity score and the action value function (Q-function) for future policies, while using cross-fitting. The Q-functions are estimated via fitted Q-evaluation \citep{Munos_et_al_2008, Fonteneau_et_al_2013, Le_et_al_2019}, a method for offline policy evaluation in reinforcement learning. The optimal policy at each stage is then estimated by maximizing the estimated policy value function over a pre-specified class of stage-specific policies. This procedure yields the estimated DTR as a sequence of policies across all stages.



The proposed approach is computationally efficient due to its stepwise backward optimization, which is particularly advantageous given the complexity of optimizing DTRs over multiple stages. Furthermore, leveraging a classification-based framework \citep{Zhao_et_al_2012, Kitagawa_Tetenov_2018a}, our approach can enhance the interpretability of DTRs by allowing for the use of interpretable policy classes, such as decision trees, for each stage. It can also accommodate various dynamic treatment problems, including optimal stopping/starting problems, by appropriately constraining the class of feasible DTRs.
An important example of an optimal stopping problem in economics is unemployment insurance programs that reduce benefit levels after a certain duration \citep[e.g.,][]{Meyer_1995,Kolsrud_et_al_2018}. In this context, a DTR determines the timing of benefit reductions for each unemployed individual.


We study the statistical properties of the proposed approach in terms of welfare regret, defined as the average outcome loss of the estimated DTR relative to the optimal one. This is a nontrivial task, because the stage-specific policies within the DTR are estimated sequentially, rather than simultaneously, and state variables are influenced by past treatments. The main theoretical contribution of this paper is to establish the convergence rate for welfare regret, linking it to the convergence rates of the nuisance component estimators (propensity scores and Q-functions) in terms of the mean-squared error (MSE). Our key result identifies conditions on the nuisance component estimators and the class of DTRs under which the resulting DTR achieves the minimax optimal convergence rate of $n^{-1/2}$ for regret.
For instance, if all nuisance components are estimated with an MSE convergence rate of $n^{-1/4}$ -- a rate attainable by many machine learning methods under suitable structural assumptions -- and the complexity of the DTR class is appropriately constrained, the resulting DTR achieves regret convergence to zero at the optimal rate of $n^{-1/2}$. This result parallels those of \cite{Athey_Wager_2020} and \cite{zhou2022offline}, who study doubly robust policy learning in single-stage settings, and aligns with the principles of double machine learning \citep{chernozhukov_et_al_2018}.

We illustrate the proposed method through an empirical application to data from Project STAR \citep[e.g.,][]{krueger_1999}. Specifically, we learn the optimal DTR for allocating each student to either a regular-size class with a teacher aide or a small-size class without one during their early education (kindergarten and grade 1). The estimated DTR uses students’ intermediate academic performance to determine the optimal class type for the subsequent grade. Our empirical results demonstrate that the optimal DTR leads to better academic outcomes for students compared to uniform class-type allocations.

%%%%%%%%%%

\subsection*{Related Literature \label{sec:related literature}}

Although many studies have explored the statistical decision/learning of treatment choice, most have focused on the single-stage setting.\footnote{A partial list includes \cite{Manski_2004}, \cite{Hirano_Porter_2009}, \citeauthor{Stoye_2009} (\citeyear{Stoye_2009,Stoye_2012}), \cite{Qian_Murphy_2011}, \cite{Bhattacharya_Dupas_2012}, \cite{Tetenov_2012}, \cite{Zhao_et_al_2012}, \cite{Kitagawa_Tetenov_2018a}, \cite{Athey_Wager_2020}, \cite{Mbakop_Tabord-Meehan_2019}, \cite{kitagawa_et_al_2021}, \cite{zhou2022offline}, and \cite{viviano_2024}, among others.} 
Among these, this study is most closely related to \cite{Athey_Wager_2020} and \cite{zhou2022offline}, who develop doubly robust policy learning in single-stage settings and show that the $n^{-1/2}$-upper bound on regret can be achieved even in observational data settings. This paper seeks to extend their approach and results to the multi-stage dynamic treatment choice problem.

There is an expanding literature on the estimation of optimal DTRs, with various methods proposed having been proposed.\footnote{\cite{Chakraborty_Murphy_2014}, \cite{Laber_et_al_2014}, \cite{Kosorok_et_al_2019}, and \cite{Li_et_al_2023} provide reviews of the literature.} Offline Q-learning \citep{Watkins_Dayan_1992} is arguably the most widely used method for estimating optimal DTRs \citep[e.g.,][]{Murphy_2005,Moodie_et_al_2012,Zhang_et_al_2018}. \cite{Murphy_2005} shows that the performance of the DTR obtained via Q-learning depends on the accuracy of the Q-function estimates; if the estimated Q-functions deviate from the true ones, the resulting DTR can be far from optimal.
Our approach also estimates Q-functions as part of its procedure. However, by leveraging propensity score models, it is more robust and accurate than Q-learning.

This study is also related to the classification-based, inverse probability weighting approach for estimating optimal DTRs \citep[e.g.,][]{Zhao_et_al_2015, Sakaguchi_2021}. This approach uses inverse propensity weighted outcomes to estimate the value function of a DTR and maximizes it to estimate the optimal DTR over a pre-specified class of DTRs. However, the use of inverse probability-weighted outcomes can sometimes lead to excessively high variance, which may result in suboptimal DTRs \citep{Doroudi_et_al_2018}. Our approach improves upon this method by incorporating models of Q-functions, thereby enhancing its overall performance.




Doubly robust estimation for optimal DTRs has also been proposed by \cite{Zhang_et_2013}, \cite{Wallace_Moodie_2015}, and \cite{Ertefaie_et_al_2021}. \cite{Zhang_et_2013} suggest estimating the optimal DTR by maximizing an AIPW estimator of the welfare function over an entire class of DTRs. However, this approach faces computational challenges for two reasons: (i) the nuisance components to be estimated depend on each specific DTR, and (ii) the method maximizes the estimated welfare function simultaneously over the entire class of DTRs. Our approach addresses these issues by (i) ensuring that the nuisance components depend only on the estimated policies for future stages, and (ii) estimating the optimal DTR through stage-wise backward optimization. Section \ref{sec:existing_approach} provides further details on this comparison.\footnote{\cite{zhang2018c} also propose a backward induction approach to estimate optimal DTRs based on estimators of Q-functions and propensity scores, but their method is not robust to misspecification of the Q-functions. See Footnote~\ref{fn:zhang18c} for further details.}
\cite{Wallace_Moodie_2015} develop a doubly robust estimation method based on Q-learning and G-estimation \citep{Robins_2004}. \cite{Ertefaie_et_al_2021} propose a doubly robust approach for Q-learning and investigate the statistical properties of the estimated parameters in the Q-functions. By contrast, our focus is on the statistical properties of welfare regret for the estimated DTR.

In the context of policy learning for optimal stopping/starting problems, \cite{Nie_et_al_2021} develop a doubly robust learning approach with computational feasibility and show upper bounds on the associated regret. Our framework encompasses this problem as a specific case.


This study also relates to the literature on (offline) reinforcement learning in terms of multi-stage decision problems. 
However, most works assume a Markov decision process, which this study does not rely on.
In the context of non-Markov decision processes, \cite{Jiang_Li_2016}, \cite{Thomas_et_al_2016}, and \cite{Kallus_Uehara_2020} propose doubly robust methods for evaluating DTRs, but they do not focus on optimizing DTRs.

Finally, in the econometric literature on dynamic treatment analysis, \cite{Heckman_Navarro_2007} and \cite{Heckman_et_al_2016} use exclusion restrictions to identify average dynamic treatment effects, though their focus is not on the identification of optimal DTRs. \cite{Han_2023} proposes a method to characterize the sharp partial ordering of counterfactual welfares of DTRs within an instrumental variable setting. \cite{Ida_et_al_2024} empirically demonstrate that the optimal DTR outperforms optimal static targeting policies in the context of energy-saving rebate programs.

\subsection*{Structure of the Paper }
The remainder of the paper is organized as follows. Section \ref{sec:setup} outlines the dynamic treatment choice problem. Section \ref{sec:doubly_robust_policy_learning} presents the doubly robust approach for learning the optimal DTRs through backward induction. Section \ref{sec:statistical_properties} shows the statistical properties of the proposed approach. Section \ref{sec:existing_approach} compares the proposed approach with that of \cite{Zhang_et_2013}. Section \ref{sec:simulation} presents a simulation study to evaluate the finite sample performance of the approach. Section \ref{sec:empirical_illustration} shows the empirical application results. All proofs of the main theorem and auxiliary lemmas are provided in Appendix.

%%%%%%%%%%%%%%%%%%%%%%%%%%%%%%%%%%%%%%%

\section{Setup\label{sec:setup}}

Section \ref{sec:dynamic treatment framework} introduces the dynamic treatment framework, following the dynamic potential outcome framework of \citet{Robins_1986,Robins_1997} and \cite{Murphy_2003}. Subsequently, we define the dynamic treatment choice problem in Section \ref{sec:dynamic treatment choice problem}. 

\subsection{Dynamic Treatment Framework \label{sec:dynamic treatment framework}}

We consider a fixed number of stages, $T$ $(< \infty)$, for multiple treatment assignments.
Let $\MA_{t} \equiv \{0,\ldots,d_t -1\}$ for $t=1,\ldots,T$ denote the set of possible treatment arms at stage $t$, where $d_t$ ($\geq 2$) is the number of treatment arms at stage $t$, which may vary across stages. We observe the assigned treatment $A_{t} \in \MA_{t}$ for each individual at each stage $t$. Let $S_{t}$ be a vector of the state variables observed prior to treatment assignment at stage $t$, which may depend on past treatments. 
At each stage $t$, we observe the outcome $Y_{t}$ after the treatment intervention. The state vector $S_{t}$ (for $t \geq 2$) may include previous outcomes $(Y_{1},\ldots,Y_{t-1})$. 
Throughout this paper, for any time-dependent object $V_{t}$, we denote by $\text{\ensuremath{\text{\ensuremath{\underline{V}}}}}_{t}\equiv \left(V_{1},\ldots,V_{t}\right)$
its history up to stage $t$, and by $\text{\ensuremath{\underline{V}}}_{s:t} \equiv \left(V_{s},\ldots,V_{t}\right)$ the partial history from stage $s$ up to stage $t$ (for $s\leq t$).
For example, $\underline{A}_{t}$ denotes the treatment history up to stage $t$.

Let $Z\equiv(\underline{A}_{T},\underline{S}_T,\underline{Y}_{T})$ denote the vector of all observed variables. We define the history at stage $t$ as $H_t \equiv (\underline{A}_{t-1},\underline{S}_t)$, which is the information available to the decision-maker when selecting the treatment at stage $t$. Note that $H_1 = (S_1)$, where $S_1$ represents individual characteristics observed prior to the beginning of the sequential treatment intervention. We denote the supports of $H_{t}$ and $Z$ by ${\cal H}_{t}$ and $\MZ$, respectively.

To formalize our results, we adopt the dynamic potential outcomes framework \citep{Robins_1986, Hernan_et_al_2001, Murphy_2003}. Let $\underline{\MA}_t \equiv \MA_1 \times \cdots \times \MA_t$. For each $\underline{a}_t \in \underline{\MA}_t$, we define $Y_{t}\left(\text{\ensuremath{\underline{a}}}_{t}\right)$ as the potential outcome for stage $t$, representing the outcome that would be realized at stage $t$ if the treatment history up to that stage were $\underline{a}_{t}$. We assume that the outcome at each stage is not influenced by treatments in future stages. Since the state variables $S_t$ may also depend on past treatments, we define the potential state variables as $S_{t}(\underline{a}_{t-1})$ for each $t\geq 2$ and $\underline{a}_{t-1} \in \underline{\MA}_{t-1}$. For $t = 1$, we set $S_{1}(\underline{a}_{0}) = S_1$ for notational convenience. The observed outcomes and state variables are thus defined as $Y_{t} \equiv Y_{t}(\underline{A}_t)$ and $S_t \equiv S_t(\underline{A}_{t-1})$, respectively. 

Let $\underline{S}_t(\underline{a}_{t-1}) \equiv (S_1,S_2(\underline{a}_{1}),\ldots,S_t(\underline{a}_{t-1}))$ be the history of the potential state variables. We define the vector $H_{t}\left(\underline{a}_{t-1}\right) \equiv \left(\underline{a}_{t-1},\underline{S}_t(\underline{a}_{t-1})\right)$  as the potential history realized when the prior treatments are $\underline{a}_{t-1}$. For $t=1$, we set $H_{1}\left(\underline{a}_{0}\right) = H_1$. The observed history is then defined as $H_t \equiv H_t\left(\underline{A}_{t-1}\right)$.
We denote by $P$ the distribution of all underlying variables $\left(\underline{A}_T,\left\{\underline{S}_T(\underline{a}_{T-1})\right\}_{\underline{a}_{T-1} \in \underline{\MA}_{T-1}},\{\underline{Y}_{T}(\underline{a}_T)\}_{\underline{a}_T \in \underline{\MA}_T}\right)$, where  $\underline{Y}_T(\underline{a}_T) \equiv (Y_1(a_1),Y_2(\underline{a}_{2}),\ldots,Y_T(\underline{a}_{T}))$.

From an observational study, we observe $Z_{i}\equiv \left(\underline{A}_{i,T},\underline{S}_{i,T},\underline{Y}_{i,T}\right)$ for individuals $i=1,\ldots,n$, where $\underline{A}_{i,T}=(A_{i,1},\ldots,A_{i,T})$, $\underline{S}_{i,T}=(S_{i,1},\ldots,S_{i,T})$, and $\underline{Y}_{i,T}=(Y_{i,1},\ldots,Y_{i,T})$. The observed outcome $Y_{i,t}$ and state variables $S_{i,t}$ are defined as $Y_{i,t}  \equiv Y_{i,t}(\underline{A}_{i,t})$ and $S_{i,t}  \equiv S_{i,t}(\underline{A}_{i,t-1})$, respectively, with $Y_{i,t}\left(\underline{a}_{t}\right)$ and $S_{i,t}\left(\underline{a}_{t-1}\right)$ being the
potential outcome and state variables, respectively. Let $\underline{S}_{i,T}(\underline{a}_{T-1}) \equiv (S_{i,1},S_{i,2}(a_1),\ldots,S_{i,T}(\underline{a}_{T-1}))$ and $\underline{Y}_{i,T}(\underline{a}_{T}) \equiv (Y_{i,1}(a_1),\ldots,Y_{i,T}(\underline{a}_{T}))$. We assume that the vectors of random variables $\left(\underline{A}_{i,T},\left\{\underline{S}_{i,T}\left(\text{\ensuremath{\underline{a}}}_{T-1}\right)\right\}_{\text{\ensuremath{\underline{a}}}_{T-1}\in\underline{\MA}_{T-1}}, \left\{ \underline{Y}_{i,T}\left(\text{\ensuremath{\underline{a}}}_{T}\right)\right\} _{\text{\ensuremath{\underline{a}}}_{T}\in\underline{\MA}_{T}}\right)$, for $i=1,\ldots,n$, are independent and identically distributed (i.i.d) under the distribution $P$. 
We denote by $H_{i,t}\equiv (\underline{A}_{i,t-1},\underline{S}_{i,t})$ the history of the $i$-th individual at stage $t$.

We define $e_{t}\left(h_{t},a_{t}\right)\equiv \BP\left(A_{t}=a_{t}\mid H_{t}=h_{t}\right)$ as the propensity score of treatment $a_t$ at stage $t$ with  history $h_t$. In the observational data setting we study, the propensity scores are unknown to the analyst. This contrasts with the experimental data setting, in which the propensity scores are known from the experimental design.\footnote{Even when data are obtained from a sequential multiple assignment randomized trial (SMART), propensity scores may be unknown due to non-compliance with assigned treatments or attrition at certain stages.}

Throughout the paper, we assume that the underlying distribution $P$ satisfies the following assumptions.

\medskip

\begin{assumption}[Sequential Ignorability]\label{asm:sequential independence} For any $t=1,\ldots,T$ and $\text{\ensuremath{\underline{a}}}_{T}\in\underline{\MA}_{T}$,
\begin{align*}
\{Y_{t}\left(\underline{a}_{t}\right),\ldots,Y_{T}\left(\underline{a}_{T}\right),S_{t+1}(\underline{a}_{t}),\ldots,S_{T}(\underline{a}_{T-1})\} \indep A_{t} \mid H_{t}.  
\end{align*}
\end{assumption}

%\smallskip

\begin{assumption}[Bounded Outcomes]\label{asm:bounded outcome}
There exists $M <\infty$ such that the support of $Y_{t}(\underline{a}_{t})$ is
contained in $\left[-M/2,M/2\right]$ for all $t \in \{1,\ldots,T\}$ and $\underline{a}_{t} \in \underline{\MA}_{t}$.
\end{assumption}

%\smallskip

\medskip


Assumption \ref{asm:sequential independence} is also referred to as the dynamic unconfoundedness assumption or sequential conditional independence assumption, and is commonly used in the literature on dynamic treatment effect analysis \citep{Robins_1997,Murphy_2003}. This assumption implies that the treatment assignment at each stage is independent of current and future potential outcomes, as well as future state variables, conditional on the history up to that point. In observational studies, this assumption holds when a sufficient set of confounders is controlled for at each stage. Assumptions \ref{asm:bounded outcome} is a standard assumption in the literature on estimating optimal DTRs.

%\medskip

\subsection{Dynamic Treatment Choice Problem\label{sec:dynamic treatment choice problem}}

The aim of this study is  to develop a method for learning optimal DTRs using data from an observational study. We define a policy for each stage $t$ as $\pi_{t}:\MH_{t}\mapsto \MA_{t} $, a map from the history space for stage $t$ to the treatment space for stage $t$. A policy $\pi_t$ determines which treatment is assigned to each individual at stage $t$ based on their history $h_t$. We define a DTR as $\pi \equiv (\pi_{1}, \ldots, \pi_{T})$, a sequence of stage-specific policies. The DTR guides the treatment choice for each individual from the first to the final stage, based on their evolving history up to each stage.

Given a fixed DTR $\pi$, we define the welfare of $\pi$ as
\begin{align*}
W\left(\pi\right) &\equiv   
\E\left[\sum_{t=1}^{T}\sum_{\underline{a}_t \in \underline{\MA}_t}\left( Y_{t}(\underline{a}_t)\cdot \prod_{s=1}^{t}\mathbf{1}\{\pi_s(H_{s}(\underline{a}_{s-1}))=a_s\}\right)\right].
\end{align*}
This expression represents the expected total outcome realized when treatments are assigned sequentially according to the DTR $\pi$.\footnote{We can also define and consider welfare with weighted outcomes as follows: $W\left(\pi\right) =   
\E\left[\sum_{t=1}^{T}\sum_{\underline{a}_t \in \underline{\MA}_t}\left(\gamma_{t} Y_{t}(\underline{a}_t)\cdot \prod_{s=1}^{t}\mathbf{1}\{\pi_s(H_{s}(\underline{a}_{s-1}))=a_s\}\right)\right]$, where $\gamma_{t}$ represents the weight assigned to the outcome at stage $t$. For instance, $\gamma_t$ can be a discount factor $\gamma^t$ with $\gamma$ being a discount rate. If we focus solely on the final-stage outcome, we set $\gamma_{T}=1$ and $\gamma_{1}=\cdots=\gamma_{T-1}=0$.}

We consider choosing a DTR from a pre-specified class of DTRs denoted by $\Pi \equiv \Pi_{1}\times\cdots\times \Pi_{T}$, where $\Pi_{t}$ represents a class of policies for stage $t$ (i.e., a class of measurable functions $\pi_{t}:\MH_{t}\rightarrow \MA_{t}$). For example, \cite{Laber_Zhao_2015}, \cite{Tao_et_al_2018}, \cite{Sun_Wang_2021},  \cite{Blumlein_et_al_2022}, and \cite{zhou_et_al_2023} use a class of decision trees \citep{Breiman84book} for $\Pi_t$, and \cite{Zhang_et_al_2018} use a class of list-form policies for $\Pi_t$. These policy classes enhance the interpretability of the resulting DTRs.%\footnote{\cite{Ahmad_et_al_2018} discuss the importance of interpretability of machine learning models in healthcare.} 

Our framework also accommodates cases where only a part of the history, $\tilde{h}_t$, a subvector of $h_{t}$, is used for treatment choice by constraining $\Pi_t$ to be a class of functions of $\tilde{h}_t$. This is particularly relevant given the increasing dimension of the entire history $h_t$ over time, where only a sub-history may be informative for optimal treatment choice. 
Additionally, our framework can encompass the optimal stopping/starting problem by constraining the class $\Pi$ of DTRs as follows.

\medskip

\begin{example}[Optimal Starting/Stopping Problem]\label{eg:optimal_start}
Suppose that the number of treatment arms $d_t$ is time-invariant ($d_t = K$ for some constant $K$) and that arm $0$ represents no treatment. Our framework can accommodate the optimal starting problem, in which the decision-maker determines when to start assigning one of the arms $a_t \in \{1,\ldots,K-1\}$ for each unit. This problem can be incorporated into our framework by constraining the class $\Pi$ of DTRs such that for any $t \in \{2,\ldots,T\}$ and $(\pi_{t}, h_{t}) \in \Pi_{t} \times \MH_{t}$, $\pi_{t}(h_t) = a_{t-1}$ if $a_{t-1} \neq 0$. The optimal stopping problem can also be specified in a similar manner.
\end{example}
\medskip

An important example of an optimal stopping problem in economics is unemployment insurance programs that reduce benefit levels after a certain duration \citep[e.g.,][]{Meyer_1995,Kolsrud_et_al_2018}. In this context, a DTR determines the timing of benefit reductions for each unemployed individual.

Given a pre-specified class $\Pi$ of DTRs, we impose an overlap condition on $\underline{A}_{T}$, related to the structure of $\Pi$, as follows.

\medskip

\begin{assumption}[Overlap Condition]\label{asm:overlap} 
There exists $\eta \in (0,1)$ for which $\eta \leq e_{t}(h_{t},a_{t})$ holds for any $t \in \{1,\ldots,T\}$ and any pair  $(h_{t},a_{t}) \in \MH_{t} \times \MA_{t}$ such that there exists $\pi_{t} \in \Pi_{t}$ that satisfies $\pi_{t}(h_t) = a_{t}$. 
\end{assumption}

\medskip

When $\Pi$ is structurally constrained, Assumption \ref{asm:overlap} is weaker than the common overlap condition that requires $e_{t}(h_t, a_t) \in (0, 1)$ for all $(h_t, a_t) \in \MH_{t} \times \MA_{t}$ and $t$. For example, in the optimal stopping problem, Assumption  \ref{asm:overlap} does not require $e_{t}(h_t, 0) > 0$ for any $h_t$ such that $a_s$ in $h_t$ is 0 for some $s < t$.

The ultimate goal of our analysis is to choose an optimal DTR that maximizes the welfare $W\left(\cdot\right)$ over $\Pi$. 
We are especially interested in learning the optimal DTR from observational data that satisfies the sequential ignorability assumption (Assumption \ref{asm:sequential independence}). The following section presents a novel step-wise approach to learning the optimal DTR.



\section{Learning of the Optimal DTR}\label{sec:doubly_robust_policy_learning}

In this section, we propose a doubly robust approach to learning the optimal DTRs through backward induction. Section \ref{sec:fitted_Q} first introduces Q-function (action-value function) and the fitted Q-evaluation, a method to estimate the Q-functions. This section also discusses identifiability of the optimal DTR through the backward-induction procedure. 
Section \ref{sec:learning_DTR} then presents our proposed approach to learning the optimal DTRs.

\subsection{Fitted Q-evaluation and Backward Induction}
\label{sec:fitted_Q}

For any DTR $\pi$ and class $\Pi$ of DTRs, we denote their partial sequences by $\pi_{s:t} \equiv (\pi_s,\ldots,\pi_t)$ and $\Pi_{s:t} \equiv \Pi_s \times \cdots \times \Pi_t$, respectively, for $s\leq t$.\footnote{For any object $v_{s:t}$ and $\underline{w}_{s:t}$ ($s\leq t$), $v_{t:t}$ and $\underline{w}_{t:t}$ correspond to $v_t$ and $w_t$, respectively.}
We define the policy value of $\pi_{t:T}$ for stage $t$ as 
\begin{align*}
    V_{t}\left(\pi_{t:T}\right) \equiv \E\left[\sum_{s=t}^{T}\sum_{\underline{a}_{s} \in \underline{\MA}_{s}}\left( Y_{s}(\underline{a}_{s})\cdot \mathbf{1}\{\underline{A}_{t-1} = \underline{a}_{t-1}\} \cdot \prod_{\ell=t}^{s}\mathbf{1}\{\pi_{\ell}\left(H_{\ell}(\underline{a}_{\ell-1})\right)=a_\ell\}\right)\right],
\end{align*}
where we assume $\mathbf{1}\{\underline{A}_{0} = \underline{a}_{0}\}=1$ for $t=1$. $V_{t}\left(\pi_{t:T}\right)$
represents the average total outcome from stage $t$ to stage $T$ that is realized when the treatment assignments before stage $t$ follow $\underline{A}_{t-1}$ (i.e., assignments in the observational data), and those from stage $t$ follow $\pi_{t:T}$.
With some abuse of terminology, we refer to $V_{t}\left(\pi_{t:T}\right)$ as the policy value function of $\pi_{t:T}$.
Note that $V_{1}(\pi_{1:T}) = W(\pi)$.

Given a fixed DTR $\pi$, we define Q-functions (state-action-value functions), recursively, as follows:
\begin{align}
    Q_T(h_T,a_T) &\equiv \E\left[Y_{T}|H_T = h_T, A_T = a_T\right], \label{eq:action_value_func_1}\\
     Q_{T-1}^{\pi_{T}}(h_{T-1},a_{T-1}) &\equiv \E\left[Y_{T-1} + Q_T(h_T,\pi_T(H_{T}))|H_{T-1} = h_{T-1}, A_{T-1} = a_{T-1}\right],
     \end{align}
     and, for $t=T-2,\ldots,1$,
     \begin{align}
     Q_{t}^{\pi_{(t+1):T}}(h_t,a_t) &\equiv \E\left[Y_{t} + Q_{t+1}^{\pi_{(t+2):T}}(H_{t+1},\pi_{t+1}(H_{t+1}))|H_t = h_t, A_t = a_t\right]. \label{eq:state_value_func_2}
\end{align}
We refer to $Q_{t}^{\pi_{(t+1):T}}(h_t,a_t)$ as the Q-function for $\pi_{(t+1):T}$. 
The Q-function $Q_{t}^{\pi_{(t+1):T}}(h_t,a_t)$ represents the average total outcome when the history $H_t$ and treatment $A_t$ correspond to $h_t$ and $a_t$ at stage $t$, and the future treatment assignments follow $\pi_{(t+1):T}$. 
When $t=T$, we denote $Q_{T}^{\pi_{(T+1):T}}(\cdot,\cdot) = Q_{T}(\cdot,\cdot)$.  Note that $\E[Q_{t}^{\pi_{(t+1):T}}(H_{t},\pi_{t}(H_{t}))] = V_{t}(\pi_{t:T})$ and $\E[Q_{1}^{\pi_{2:T}}(H_{1},\pi_{1}(H_{1}))] = W(\pi)$ hold under Assumption \ref{asm:sequential independence}.\footnote{See, for example, \citeauthor{Tsiatis_et_al_2019}(\citeyear{Tsiatis_et_al_2019}, Section 6.4) or Lemma \ref{lemma:identification}.} 
In what follows, for any function $f(\cdot,\cdot):\MH_{t} \times \MA_{t} \rightarrow \Real$ and policy $\pi_{t}(\cdot):\MH_{t} \rightarrow \MA_{t}$, we denote $f(h_{t},\pi_{t}(h_{t}))$ shortly by $f(h_{t},\pi_{t})$ (e.g., $Q_{t}^{\pi_{(t+1):T}}(h_t,\pi_t(h_t))$ is denoted by $Q_{t}^{\pi_{(t+1):T}}(h_t,\pi_t)$).



Given a fixed DTR $\pi$, we can use the sequential definitions in equations (\ref{eq:action_value_func_1})--(\ref{eq:state_value_func_2}) to estimate the sequence $\{Q_{t}^{\pi_{(t+1):T}}(\cdot,\cdot)\}_{t=1,\ldots,T}$ of the Q-functions for $\pi$. This approach is referred to as the fitted Q-evaluation \citep{Munos_et_al_2008,Fonteneau_et_al_2013,Le_et_al_2019} in the  reinforcement learning literature and comprises multiple steps as follows:
\begin{itemize}
    \item Regress $Y_{T}$ on $(H_T,A_T)$ to obtain $\widehat{Q}_{T}(\cdot,\cdot)$ as the estimated regression function for $Q_T(\cdot,\cdot)$;
    \item Regress $Y_{T-1} + \widehat{Q}_{T}(H_{T},\pi_{T})$ on $(H_{T-1},A_{T-1})$ to obtain $\widehat{Q}_{T-1}^{\pi_{T}}(\cdot,\cdot)$ as the estimated regression function for $Q_{T-1}^{\pi_{T}}(\cdot,\cdot)$;
   \item Recursively, for $t=T-2,\ldots,1$, regress $Y_{t} + \widehat{Q}_{t+1}^{\pi_{(t+2):T}}(H_{t+1},\pi_{t+1})$ on $(H_t,A_t)$ to obtain $\widehat{Q}_{t}^{\pi_{(t+1):T}}(\cdot,\cdot)$ as the estimated regression function for $Q_{t}^{\pi_{(t+1):T}}(\cdot,\cdot)$.
\end{itemize}
We can apply flexible/nonparametric regression methods, including machine learning methods (e.g., random forests, lasso, neural networks), to the regression in each step. 


Given the definitions of the Q-functions, we can optimize the DTR through backward induction. To present the idea, we here assume that the generative distribution function $P$ is known and that the pair  $(P,\Pi)$ satisfies Assumptions \ref{asm:sequential independence}, \ref{asm:bounded outcome}, and \ref{asm:overlap}. The backward-induction approach in the population problem is a step-wise process proceeding as follows. 
In the first step, the approach optimizes the final-stage policy over $\Pi_T$ by solving 
\begin{align*}
        \pi_{T}^{\ast,B} = \argmax_{\pi_{T} \in \Pi_{T}} \E\left[Q_{T}(H_{T},\pi_{T})\right].
\end{align*}
Then, recursively, from $t=T-1$ to $1$, the approach optimizes the $t$-th stage policy over $\Pi_{t}$ by solving
\begin{align*}
        \pi_{t}^{\ast,B} = \argmax_{\pi_{t} \in \Pi_{t}} \E\left[Q_{t}^{\pi_{(t+1):T}^{\ast,B}}(H_{t},\pi_{t})\right].
\end{align*}
Note that the objective function $\E\left[Q_{t}^{\pi_{(t+1):T}^{\ast,B}}(H_{t},\pi_{t})\right]$ corresponds to the policy value $V_{t}(\pi_{t},\pi_{(t+1):T}^{\ast,B})$ under the sequential ignorability assumption (Assumption \ref{asm:sequential independence}).
The entire procedure yields the DTR $\pi^{\ast,B} \equiv (\pi_{1}^{\ast,B},\cdots,\pi_{T}^{\ast,B})$.

We denote the optimal DTR over $\Pi$ by $\pi^{\ast,opt} = \argmax_{\pi \in \Pi} W(\pi)$. Importantly, the DTR $\pi^{\ast,B}$ obtained through backward induction does not necessarily correspond to the optimal DTR $\pi^{\ast,opt}$ when $\Pi$ is structurally constrained, as noted by \cite{Li_et_al_2023} and \cite{Sakaguchi_2021}. 
To ensure that backward induction yields the optimal DTR, the policy class $\Pi_t$ for each stage $t$ ($\geq 2$) needs to be correctly specified. The following assumption gives a sufficient condition for the backward optimization to attain optimality.

\medskip

\begin{assumption}[Correct Specification] 
\label{asm:first-best} 
There exists $\pi_{2:T}^{\ast} = (\pi_{2}^{\ast},\ldots,\pi_{T}^{\ast})\in \Pi_{2:T}$ such that for any $t=2,\ldots,T$,
\begin{align*}
Q_{t}^{\pi_{(t+1):T}^{\ast}}\left(H_t,\pi_{t}^{\ast}\right) \geq \sup_{\pi_{t} \in \Pi_t}Q_{t}^{\pi_{(t+1):T}^{\ast}}\left(H_t,\pi_{t}\right) \mbox{\ a.s.}
\end{align*} 
\end{assumption}

\medskip

Assumption \ref{asm:first-best} requires that for each stage $t \geq 2$, $\Pi_t$ includes a policy that can select the optimal arm among all feasible arms in $\{\pi_{t}(h_t):\pi_{t} \in \Pi_t\}$ $(\subseteq \MA_{t})$ for any history $h_t$. This assumption is satisfied when $\Pi_t$ (for $t=2,\ldots,T$) is flexible enough or correctly specified in relation to treatment effect heterogeneity.\footnote{ \cite{Zhang_et_2013} discuss how to correctly specify $\Pi_t$ based on the model for treatment effect heterogeneity. \cite{Zhao_et_al_2015} use a reproducing kernel Hilbert space for each $\Pi_t$ as a flexible class of polices.} We suppose that $\Pi$ satisfies Assumption \ref{asm:first-best}. 

Note that Assumption \ref{asm:first-best} is a sufficient but not a necessary condition. %To illustrate this, Appendix \ref{app:suboptimallity_backward_induction} in the supplement provides two examples demonstrating cases where backward induction both succeeds and fails to achieve optimality when Assumption \ref{asm:first-best} is not satisfied.
To illustrate this, Appendix \ref{app:suboptimallity_backward_induction} provides an example demonstrating cases where backward induction can achieve optimality even when Assumption \ref{asm:first-best} is not satisfied.

A stronger version of Assumption \ref{asm:first-best} is that each $\Pi_t$ ($t \geq 2$) contains the first-best policy; that is, there exists $\pi_{2:T}^{\ast,FB} = (\pi_{2}^{\ast,FB},\ldots,\pi_{T}^{\ast,FB})\in \Pi_{2:T}$ such that for any $t=2,\ldots,T$,
\begin{align*}
    Q_{t}^{\pi_{(t+1):T}^{\ast,FB}}\left(H_t,\pi_{t}^{\ast,FB}\right) \geq \max_{a_{t} \in \MA_t}Q_{t}^{\pi_{(t+1):T}^{\ast,FB}}\left(H_t,a_t\right) \mbox{\ a.s.}
\end{align*}
The first-best policy $\pi_{t}^{\ast,FB}$ is the policy that selects the best treatment arm for any history $h_t$. \cite{Li_et_al_2023} and \cite{Sakaguchi_2021} argue that the availability of the first-best policies is a fundamental assumption for the optimality of the backward optimization. However, Assumption \ref{asm:first-best} is practically weaker than this.
For example, in the optimal starting problem (Example \ref{eg:optimal_start}), although $\Pi_t$ may not include the first-best policy $\pi_{t}^{\ast,FB}$ due to structural constraints from the optimal starting problem, Assumption \ref{asm:first-best} does not require its feasibility. 
Similarly, when each policy in $\Pi_t$ depends only on a sub-history $\tilde{h}_t$ (a subvector of $h_t$), Assumption \ref{asm:first-best} requires that the optimal policy is available for the sub-history, rather than for the entire history.

%The following lemma shows that the backward induction procedure leads to the optimal DTR under Assumption \ref{asm:first-best}. 

The following lemma formalizes the optimality of the backward-induction procedure under Assumption \ref{asm:first-best}.

\medskip

\begin{lemma}\label{lem:optimality_backward_induction}
Under Assumptions \ref{asm:sequential independence}, \ref{asm:overlap}, and \ref{asm:first-best}, $\pi^{\ast,B}$ is the optimal DTR over $\Pi$; i.e.,
\begin{align*}
    W(\pi^{\ast,B}) \geq W(\pi) \mbox{\ for all\ }\pi \in \Pi.
\end{align*}
\end{lemma}

\begin{proof}
See Appendix \ref{app:preliminary_results}.
\end{proof}


%%%%%%%%%%%%%%%%%%%%%%%%%%%%%%%%%%%%%%%%%%%%%%%%%%%%%%%%%%%%


\subsection{Learning of the Optimal DTRs through Backward Induction}
\label{sec:learning_DTR}

This section presents a backward optimization procedure to learn the optimal DTRs using an AIPW estimator of the policy value function. Following the doubly robust policy learning approaches of \cite{Athey_Wager_2020} and \cite{zhou2022offline}, we employ cross-fitting to estimate the policy value function and learn the optimal policy independently. We randomly divide the dataset $\{Z_i:i=1,\ldots,n\}$ into $K$ evenly-sized folds (e.g., $K=5$). Let $I_k \subseteq \{1,\ldots,n\}$ be the set of indices of the data in the $k$-th fold, and $I_{-k}$ denote the set of indices of the data excluded from the $k$-th fold. In what follows, for any statistic $\hat{f}$, we denote by $\hat{f}^{-k}$ the corresponding statistic calculated using the data excluded from the $k$-th fold. Let $k(i)$ denote the fold number that contains the $i$-th observation.


The proposed approach is based on backward induction and thus consists of multiple steps. As a preliminary step, we estimate the propensity scores $\{e_t(\cdot,\cdot)\}_{t=1,\ldots,T}$ for all stages and the Q-function $Q_T(\cdot,\cdot)$ for the final stage using the data excluded in each cross-fitting fold. For each $k$, we denote by $\hat{e}_{t}^{-k}(\cdot,\cdot)$ and $\widehat{Q}_{T}^{-k}(\cdot,\cdot)$, respectively, the estimators of $e_t(\cdot,\cdot)$ and $Q_T(\cdot,\cdot)$ using data not contained in the $k$-th cross-fitting fold. Any regression methods, including machine learning methods (e.g., random forests and neural networks), can be used to estimate $e_t(\cdot,\cdot)$ and $Q_{T}(\cdot,\cdot)$. 

Given $\{\hat{e}_{t}^{-k}(\cdot,\cdot)\}_{t=1,\ldots,T}$ and  $\widehat{Q}_{T}^{-k}(\cdot,\cdot)$ for each $k=1,\ldots,K$, we estimate the optimal DTR sequentially as follows. In the first step, regarding the final stage $T$, we construct a score function of the treatment $a_T$ for stage $T$ as 
\begin{align}
    \widehat{\Gamma}_{i,T}(a_T) \equiv \frac{Y_{i,T} - \widehat{Q}_{T}^{-k(i)}(H_{i,T},A_{i,T})}{\hat{e}_{T}^{-k(i)}(H_{i,T},A_{i,T})}\cdot \mathbf{1}\{A_{i,T}=a_T\} +
    \widehat{Q}_{T}^{-k(i)}(H_{i,T},a_{T}). \label{eq:score_func_T}
\end{align}
Given a policy $\pi_T$, the sample mean $(1/n)\sum_{i=1}^{n}\widehat{\Gamma}_{i,T}(\pi_T(H_{i,T}))$ is an AIPW estimator of the policy value $V_{T}(\pi_{T})$ for stage $T$.

We then find the best candidate policy for stage $T$ by solving
\begin{align*}
\hat{\pi}_{T} = \argmax_{\pi_{T}\in \Pi_{T}}\frac{1}{n}\sum_{i=1}^{n}\widehat{\Gamma}_{i,T}\left(\pi_T(H_{i,T})\right). %\label{eq:optimization_T}
\end{align*}


In the next step, we consider stage $T-1$. Given the estimated policy $\hat{\pi}_{T}$ from the previous step, for each $k$-th cross-fitting fold, we estimate the Q-function  $Q_{T-1}^{\hat{\pi}_{T}} (\cdot,\cdot)$ for $\hat{\pi}_{T}$ by regressing $Y_{i,T-1} + \widehat{Q}_{T}(H_{i,t},\hat{\pi}_{T})$ on $(H_{i,T-1},A_{i,T-1})$ using the observations whose indices are not included in $I_k$. This corresponds to the second step of the fitted-value Q-evaluation, and any regression method can be applied at this step. We denote the resulting estimator of $Q_{T-1}^{\hat{\pi}_{T}}(\cdot,\cdot)$ by $\widehat{Q}_{T-1}^{\hat{\pi}_{T},-k}(\cdot,\cdot)$ for each $k$. 

We subsequently construct the score function of $a_{T-1}$ as follows:
\begin{align*}
    \widehat{\Gamma}_{i,T-1}^{\hat{\pi}_{T}}(a_{T-1})  &\equiv \frac{ Y_{i,T-1} + \widehat{\Gamma}_{i,T}(\hat{\pi}_{T}(H_{i,T})) - \widehat{Q}_{T-1}^{\hat{\pi}_{T},-k(i)}(H_{i,T-1},A_{i,T-1}) }{\hat{e}_{T-1}^{-k(i)}(H_{i,T-1},A_{i,T-1})} \cdot \mathbf{1}\{A_{i,T-1} = a_{T-1}\} \\
    &+ 
\widehat{Q}_{T-1}^{\hat{\pi}_{T},-k(i)}(H_{i,T-1},a_{T-1}).
\end{align*}
Given a policy $\pi_{T-1}$, the sample mean $(1/n)\sum_{i=1}^{n}\widehat{\Gamma}_{i,T-1}^{\hat{\pi}_{T}}(\pi_{T-1}(H_{i,T-1}))$ is an AIPW estimator of the policy value $V_{T-1}(\pi_{T-1},\hat{\pi}_{T})$.
We then find the best candidate policy for stage $T-1$ by solving
\begin{align*}
    \hat{\pi}_{T-1} = \argmax_{\pi_{T-1}\in \Pi_{T-1}} \frac{1}{n}\sum_{i=1}^{n}\widehat{\Gamma}_{i,T-1}^{\hat{\pi}_T}\left(\pi_{T-1}(H_{i,T-1})\right).
\end{align*}
%Given the solution $\hat{\pi}_{T-1}$, the sample mean $(1/n)\sum_{i=1}^{n}\widehat{\Gamma}_{i,T-1}^{\hat{\pi}_T}\left(\hat{\pi}_{T-1}(H_{i,T-1})\right)$ becomes an AIPW estimator of the policy value $V_{T-1}(\hat{\pi}_{(T-1):T})$ for the sequence of the estimated policies $\hat{\pi}_{(T-1):T}$.

Recursively, for $t=T-2,\ldots,1$, we learn the optimal policy as follows. For each cross-fitting index $k$, we first estimate the Q-function $Q_{t}^{\hat{\pi}_{(t+1):T}}$ by regressing  $Y_{i,t} + \widehat{Q}_{t+1}^{\hat{\pi}_{(t+2):T}}(H_{i,t+1},\hat{\pi}_{t+1})$ on $(H_{i,t},A_{i,t})$ using the observations whose indices are not in $I_k$ (the fitted Q-evaluation). Any regression method can be applied at this step. 

We next construct the score function of $a_{t}$ as
\begin{align}
    \widehat{\Gamma}_{i,t}^{\hat{\pi}_{(t+1):T}}(a_t) &\equiv \frac{ Y_{i,t} + \widehat{\Gamma}_{i,t+1}^{\hat{\pi}_{(t+2):T}}(\hat{\pi}_{t+1}(H_{i,t+1})) - \widehat{Q}_{t}^{\hat{\pi}_{(t+1):T},-k(i)}(H_{i,t},A_{i,t}) }{\hat{e}_{t}^{-k(i)}(H_{i,t},A_{i,t})}\cdot \mathbf{1}\{A_{i,t} = a_t\} \nonumber \\
    &+ 
    \widehat{Q}_{t}^{\hat{\pi}_{(t+1):T},-k(i)}(H_{i,t},a_{t}). \label{eq:score_func_t}
\end{align}
We then find the best candidate policy for stage $t$ by solving
\begin{align*}
    \hat{\pi}_{t} = \argmax_{\pi_{t}\in \Pi_{t}} \frac{1}{n}\sum_{i=1}^{n} \widehat{\Gamma}_{i,t}^{\hat{\pi}_{(t+1):T}}\left(\pi_{t}(H_{i,t})\right), %\label{eq:optimization_t}
\end{align*}
where the objective function $(1/n)\sum_{i=1}^{n} \widehat{\Gamma}_{i,t}^{\hat{\pi}_{(t+1):T}}\left(\pi_{t}(H_{i,t})\right)$ is an AIPW estimator of the policy value $V_{t}(\pi_t,\hat{\pi}_{(t+1):T})$.

Throughout this procedure, we obtain the sequence $\hat{\pi} \equiv (\hat{\pi}_{1},\ldots,\hat{\pi}_{T})$, which serves as the estimator for the optimal DTR. Algorithm \ref{alg:proposed} summarizes the entire procedure.\footnote{\label{fn:zhang18c} Our proposed method differs from that of \cite{zhang2018c} in the construction of the objective function $\widehat{\Gamma}_{i,t}^{\hat{\pi}{(t+1):T}}$, even in the binary treatment setting considered by \cite{zhang2018c}, although both approaches use Q-functions and propensity scores. In their approach, the objective function at each stage $t$ depends only on the propensity score for that stage and the Q-functions for that and future stages. As a result, it is not robust to misspecification of Q-functions for future stages -- even when the propensity scores are correctly specified -- and thus does not attain double robustness. Specifically, in their framework (following the notation in \cite{zhang2018c}), if $Q_k$ is misspecified, then $\widetilde{V}_k$ is no longer a consistent estimator of $V_k$, which in turn leads to inconsistency in their estimators of the policy value, $\widehat{C}_k$, and the optimal policy, $\hat{g}_{C,k}^{\mathrm{opt}}$.} In the following section, we will show statistical properties of the resulting DTR $\hat{\pi}$ with respect to its welfare regret. \\


%%%%%%%%%%%%%%%%%%%%%%%%%%%%%%%

%\begin{comment}
\begin{algorithm}[h]
\caption{Doubly robust backward-induction learning}\label{alg:proposed}
\begin{algorithmic}[1]
\State \textbf{Input:} $K$ (the number of cross-fitting folds); $\{Z_{i}:i=1,\ldots,n\}$ (dataset) %\Comment{e.g., $K=5$}
\State Construct $\hat{e}_{t}^{-k}(\cdot,\cdot)$ ($t=1,\ldots,T$) and $\widehat{Q}_{T}^{-k}(\cdot,\cdot)$ for each $k \in \{1,\ldots,K\}$
\State Compute $\widehat{\Gamma}_{i,T}(a_{T})$ (equation (\ref{eq:score_func_T})) for each $a_{T} \in \MA_{T}$ and $i \in \{1,\ldots,n\}$
\State Learn the optimal policy by setting $\hat{\pi}_{T} = \argmax_{\pi_{T}\in \Pi_{T}}(1/n)\sum_{i=1}^{n}\widehat{\Gamma}_{i,T}\left(\pi_T(H_{i,T})\right)$ %(\ref{eq:optimization_T}) %\Comment{Policy optimization}
\For{$t = T-1$ \textbf{to} $1$}
    \State Construct $\widehat{Q}_{t}^{\hat{\pi}_{(t+1):T},-k} (\cdot,\cdot)$ for each $k \in \{1,\ldots,K\}$ (Fitted Q-evaluation)
    \State Compute $\widehat{\Gamma}^{\hat{\pi}_{(t+1):T}}_{i,t}(a_{t})$ (equation (\ref{eq:score_func_t})) for each $a_{t} \in \MA_{t}$ and $i \in \{1,\ldots,n\}$
    \State Learn the optimal policy by setting $$\hat{\pi}_{t} = \argmax_{\pi_{t}\in \Pi_{t}} \frac{1}{n}\sum_{i=1}^{n} \widehat{\Gamma}_{i,t}^{\hat{\pi}_{(t+1):T}}\left(\pi_{t}(H_{i,t})\right)$$%(\ref{eq:optimization_t}) %\Comment{Policy optimization}
\EndFor
\State \textbf{return} $\hat{\pi}=(\hat{\pi}_{1},\ldots,\hat{\pi}_{T})$
\end{algorithmic}
\end{algorithm}
%\end{comment}
\bigskip





%%%%%%%%%%%%%%%%%%%%%%%%%%%%%%%
\section{Statistical Properties}\label{sec:statistical_properties}

Given a DTR $\pi \in \Pi$, we define the regret of $\pi$ by $R(\pi) \equiv \max_{\tilde{\pi} \in \Pi}W(\tilde{\pi}) - W(\pi)$, the loss of the welfare of $\pi$ relative to the maximum welfare achievable in $\Pi$.
We study statistical properties of $\hat{\pi}$ with respect to its regret $R(\hat{\pi})$. This section shows the rate of convergence of $R(\hat{\pi})$ depending on the convergence rates of the estimators of the nuisance components, $\{\hat{e}_{t}^{-k}(\cdot,\cdot)\}_{t=1,\ldots,T}$ and  $\{\widehat{Q}_{t}^{\pi_{(t+1):T},-k}(\cdot,\cdot)\}_{t=1,\ldots,T}$, and the complexity of $\Pi$. 

Let $\widehat{Q}_{t}^{\pi_{(t+1):T},(n)}(\cdot,\cdot)$ and $\hat{e}_{t}^{(n)}(\cdot,\cdot)$ denote the estimators of the Q-function $Q_{t}^{\pi_{(t+1):T}}(\cdot,\cdot)$ for $\pi_{(t+1):T}$ and the propensity score $e_t(\cdot,\cdot)$, respectively, using a sample of size $n$ randomly drawn from the distribution $P$. For $t=T$, we denote $\widehat{Q}_{T}^{\pi_{(T+1):T},(n)}(\cdot,\cdot) = \widehat{Q}_{T}^{(n)}(\cdot,\cdot)$ for notational convenience.
We suppose that $\{\widehat{Q}_{t}^{\pi_{(t+1):T},(n)}(\cdot,\cdot)\}_{t=1,\ldots,T}$ and $\{\hat{e}_{t}^{(n)}(\cdot,\cdot)\}_{t=1,\ldots,T}$ satisfy the following assumption.

\medskip{}

\begin{assumption}[Double Robustness]\label{asm:rate_of_convergence_backward}
(i) There exists $\tau >0$ such that the following holds: For all $t=1,\ldots,T$, $s=1,\ldots,t$, and $m \in \{0,1\}$,
\begin{align*}
    \sup_{\underline{a}_{s:t} \in \underline{\MA}_{s:t}} & \E\left[ \sup_{\pi_{(t+1):T} \in \Pi_{(t+1):T}}\left(\widehat{Q}_{t}^{\pi_{(t+1):T},(n)}(H_{t},a_{t}) - Q_{t}^{\pi_{(t+1):T}}(H_{t},a_{t})\right)^{2}\right] \\
    &\times \E\left[\left(\frac{1}{\prod_{\ell=s}^{t-m}\hat{e}_{\ell}^{(n)}(H_{\ell},a_{\ell})} - \frac{1}{\prod_{\ell=s}^{t-m}e_{\ell}(H_{\ell},a_{\ell})} \right)^{2}\right] = \frac{O(1)}{n^{\tau}}.
\end{align*}
(ii) There exists $n_0 \in \Natural$ such that for any $n \geq n_0$ and $t=1,\ldots,T$, 
\begin{align*}
     \sup_{a_{t} \in \MA_{t}, \pi_{(t+1):T} \in \Pi_{(t+1):T}}\widehat{Q}_{t}^{\pi_{(t+1):T},(n)}(H_{t},a_{t}) < \infty \mbox{\ \ and\ \ }
     \min_{a_t \in \MA_t}\hat{e}_{t}^{(n)}(H_{t},a_{t}) > 0 %\mbox{\ \ hold a.s.}
\end{align*}
hold a.s.
\end{assumption}

\medskip{}


As we will see later, the $\sqrt{n}$-consistency of the regret  $R(\hat{\pi})$ to zero can be achieved when Assumption \ref{asm:rate_of_convergence_backward} (i) holds with $\tau=1$. This is not very strong or restrictive. For example, Assumption \ref{asm:rate_of_convergence_backward} (i) is satisfied with $\tau = 1$ when
\begin{align*}
&\sup_{a_{t} \in \MA_{t}}  \E\left[\sup_{\pi_{(t+1):T} \in \Pi_{(t+1):T}}\left(\widehat{Q}_{t}^{\pi_{(t+1):T},(n)}(H_{t},a_{t}) - Q_{t}^{\pi_{(t+1):T}}(H_{t},a_{t})\right)^{2}\right] = \frac{O(1)}{\sqrt{n}}\mbox{\ and\ }\\
    &\sup_{\underline{a}_{s:t} \in \underline{\MA}_{s:t}} \E\left[\left(\frac{1}{\prod_{\ell=s}^{t}\hat{e}_{\ell}^{(n)}(H_{\ell},a_{\ell})} - \frac{1}{\prod_{\ell=s}^{t}e_{\ell}(H_{\ell},a_{\ell})} \right)^{2}\right] = \frac{O(1)}{\sqrt{n}}
\end{align*}
hold for all $t=1,\ldots,T$ and $s=1,\ldots,t$. The uniform MSE convergence rate of $\widehat{Q}_{t}^{\pi_{(t+1):T},(n)}(H_{t},a_{t})$ for the fitted Q-evaluation is not a standard result, but some existing results are applicable with some modifications \citep[e.g.,][]{Zhang_et_al_2018,Kallus_Uehara_2020}.
For example, \cite{Kallus_Uehara_2020} argue that the FQE estimators $\{\widehat{Q}_{t}^{\pi_{(t+1):T},(n)}:t=1,\ldots,T\}$ for a fixed DTR $\pi$ can be viewed as M-estimators that minimize $(1/n)\sum_{t=1}^{T}\sum_{i=1}^{n}\left(Y_{i,t} + Q_{t+1}^{\pi_{(t+2):T}}(H_{i,t+1},\pi_{t+1}) - Q_{t}^{\pi_{(t+1):T}}(H_{i,t},\pi_{t}) \right)^2$ over (semi/non)parametric classes of Q-functions, where $Q_{T+1}^{\pi_{(T+2):T}}(H_{i,T+1},\pi_{T+1})=0$ by convention. This formulation enables the use of existing results on (semi/nonparametric) M-estimation to derive uniform convergence rates for $\widehat{Q}_{t}^{\pi{(t+1):T},(n)}$ for each $t$.

Note also that Assumption \ref{asm:rate_of_convergence_backward} (i) encompasses the property of double robustness; that is, Assumption \ref{asm:rate_of_convergence_backward} (i) holds if either $\widehat{Q}_{t}^{\pi_{(t+1):T},(n)}(\cdot,\cdot)$ is uniformly consistent or $\hat{e}_{t}^{(n)}(\cdot,\cdot)$ is consistent, provided that Assumption \ref{asm:rate_of_convergence_backward} (ii) holds.


We next consider the complexity of the class $\Pi$ of DTRs and the classes $\Pi_t$ of stage-specific policies. Following \cite{zhou2022offline}, we use the $\epsilon$-Hamming covering number to measure the complexity of the class of policy sequences $\pi_{s:t} \in \Pi_{s}\times \cdots \times \Pi_{t}$ for each $s$ and $t$ such that $s \leq t$. 

\medskip{}

\begin{definition}
(i) For any stages $s$ and $t$ such that $s \leq t$, given a set of history points $\{h_{t}^{(1)},\ldots,h_{t}^{(n)}\} \subseteq \MH_{t}$, we define the Hamming distance between two sequences of policies  $\pi_{s:t}, \pi_{s:t}^{\prime} \in \Pi_{s:t}$ as $d_{h}(\pi_{s:t},\pi_{s:t}^{\prime}) \equiv n^{-1}\sum_{i=1}^{n}\mathbf{1}\{\pi_s(h_{s}^{(i)})\neq \pi_{s}^{\prime}(h_{s}^{(i)}) \vee \cdots \vee \pi_t(h_{t}^{(i)})\neq \pi_{t}^{\prime}(h_{t}^{(i)})\}$,
where we denote $h_{t}^{(i)} = (\underline{a}_{t-1}^{(i)},\underline{s}_{t}^{(i)})$ and $h_{\ell}^{(i)}$ ($\ell =s,\ldots,t-1$) is the subvector of $h_{t}^{(i)}$ such that $h_{\ell}^{(i)} = (\underline{a}_{\ell-1}^{(i)},\underline{s}_{\ell}^{(i)}) \in \MH_{\ell}$. 
Let $N_{d_h}\left(\epsilon,\Pi_{s:t},\left\{h_{t}^{(1)},\ldots,h_{t}^{(n)}\right\}\right)$ be the smallest number of sequences of policies $\pi_{s:t}^{(1)},\pi_{s:t}^{(2)},\ldots$ in $\Pi_{s:t}$ such that for any $\pi_{s:t} \in \Pi_{s:t}$, there exists $\pi_{s:t}^{(i)}$ satisfying $d_h(\pi_{s:t},\pi_{s:t}^{(i)}) \leq \epsilon$.
We define the $\epsilon$-Hamming covering number of $\Pi_{s:t}$ as
\begin{align*}
    N_{d_h}(\epsilon,\Pi_{s:t})  \equiv  \sup\left\{N_{d_h}\left(\epsilon,\Pi_{s:t},\left\{h_{t}^{(1)},\ldots,h_{t}^{(n)}\right\}\right)\middle| n \geq 1, h_{t}^{(1)},\ldots,h_{t}^{(n)} \in \MH_t\right\}.
\end{align*}
(ii) We define the entropy integral of $\Pi_{s:t}$ as $\kappa(\Pi_{s:t})= \int_{0}^{1}\sqrt{\log N_{d_h}\left(\epsilon^2,\Pi_{s:t}\right)}d\epsilon$.
\end{definition}

\medskip{}

Note that when $s=t$, $N_{d_h}(\epsilon,\Pi_{s:t}) = N_{d_h}(\epsilon,\Pi_t)$ and $\kappa(\Pi_{s:t}) = \kappa(\Pi_t)$. 
We assume that the policy class $\Pi_t$ for each $t$ is not excessively complex in terms of the covering number.
\medskip{}

\begin{assumption}[Complexity of $\Pi_t$]\label{asm:bounded entropy}
For all $t=1,\ldots,T$, $N_{d_h}(\epsilon,\Pi_{t}) \leq C\exp(D(1/\epsilon)^{\omega})$ holds for any $\epsilon>0$ and some constants $C,D>0$ and $0<\omega<0.5$.
\end{assumption}

\medskip{}

This assumption implies that the covering number of $\Pi_t$ does not grow too quickly, but allows $\log N_{d_h}(\epsilon,\Pi_{t})$ to grow at a rate of  $1/\epsilon$.
This assumption is satisfied, for example, by a class of finite-depth trees (see \citeauthor{zhou2022offline} (\citeyear{zhou2022offline}, Lemma 4)). In the case of a binary action set (i.e., $|\MA_t|=2$), a VC-class of $\pi_{t}$ also satisfies Assumption \ref{asm:bounded entropy}. 
\citeauthor{zhou2022offline} (\citeyear{zhou2022offline}, Remark 4) shows that the entropy integral $\kappa(\Pi_t)$ is finite under Assumption \ref{asm:bounded entropy}. 

Regarding the class $\Pi$ of entire DTRs, its entropy integral $\kappa(\Pi)$ is finite as well under Assumption \ref{asm:bounded entropy}. 

\medskip

\begin{lemma}\label{lem:entropy_integral_bound}
Under Assumption \ref{asm:bounded entropy}, $\kappa(\Pi) < \infty$.
\end{lemma}

\begin{proof}
See Appendix \ref{app:preliminary_results}.
\end{proof}

\medskip

%\cite{zhou2022offline} give some other discussion for the covering number. For example, for the policy class of finite-depth trees (\cite{Breiman_2017}), $\log N_{d_h}(\epsilon,\Pi_{t})$ is only $O(\log(1/\epsilon))$ (see \citeauthor{zhou2022offline} (\citeyearpar{zhou2022offline}, Lemma 4)). 

The following theorem presents the main result of this paper, showing the rate of convergence for the regret of the DTR $\hat{\pi}$ obtained using the proposed approach.

\medskip

\begin{theorem}\label{thm:main_theorem_backward}
Under Assumptions \ref{asm:sequential independence}--\ref{asm:overlap}, \ref{asm:first-best}, \ref{asm:rate_of_convergence_backward}, and \ref{asm:bounded entropy},
\begin{align*}
    R(\hat{\pi}) = O_{p}\left(\kappa(\Pi) \cdot n^{-1/2}\right) + O_{p} (n^{-\min\{1/2,\tau/2\}}). 
\end{align*}
\end{theorem}

\begin{proof}
See Appendix \ref{app:main_proof}.
\end{proof}

\medskip

This theorem establishes the convergence rate of the regret $R(\hat{\pi})$ for the proposed method. When Assumption \ref{asm:rate_of_convergence_backward} (i) holds with $\tau = 1$, the approach achieves the optimal rate $n^{-1/2}$ for the regret convergence.\footnote{In the case of binary treatment at each stage, \cite{Sakaguchi_2021} shows that the minimax optimal rate of convergence for the regret is $V_{1:T}\cdot n^{-1/2}$, where $V_{1:T}$ is the VC-dimension of the class of DTRs.} This result is comparable to those of \cite{Athey_Wager_2020} and \cite{zhou2022offline}, who study doubly robust policy learning in single-stage settings.\footnote{In the single-stage binary treatment setting, \cite{Athey_Wager_2020} use the specific growth rate of the entropy in a VC class and obtain a slightly stronger result compared to using the fixed entropy class.}
The asymptotic upper bound also increases with the number of time stages $T$, through the entropy integral $\kappa(\Pi)$ of the class of DTRs.

In the proof of Theorem \ref{thm:main_theorem_backward}, we consider the derivation of the asymptotic upper bound on $R(\hat{\pi})$. This is, however, a non-trivial task because the stage-specific policies in $\hat{\pi} = (\hat{\pi}_1, \ldots, \hat{\pi}_T)$ are estimated sequentially, rather than simultaneously. If the DTR were estimated simultaneously across all stages, one could adapt the theoretical analysis of \citet{zhou2022offline}. However, the sequential nature of the estimation precludes a direct application of their approach. Despite this challenge, Appendix \ref{app:main_proof} presents novel analytical techniques to evaluate the regret of the sequentially estimated DTR.

%We conclude this section with a remark comparing the proposed approach to an existing doubly-robust approach, which simultaneously optimizes a DTR across all stages.

%%%%%%%%%%%%%%%%%%%%%%%%%%%%%%%%%%%%%%%%%%%%%%%%%%%
%%%%%%%%%%%%%%%%%%%%%%%%%%%%%%%%%%%%%%%%%%%%%%%%%%%

\section{Existing Approach}\label{sec:existing_approach}

An alternative doubly robust approach for estimating the optimal DTR is maximizing an AIPW estimator of the welfare function simultaneously over the entire class of DTRs.  Specifically, we have 
\begin{align}
    \hat{\pi}^{AIPW} = \argmax_{\pi \in \Pi} & \widehat{W}^{AIPW}(\pi) \mbox{\ \ with} \label{eq:simultaneous_maximization}\\
    \widehat{W}^{AIPW}(\pi) = \frac{1}{n}\sum_{i=1}^{n}\sum_{t=1}^{T}&\left(\hat{\psi}_{it}^{-k(i)}\left(\underline{\pi}_{t}\right)\cdot Y_{it} \right. \nonumber \\
      - &\left.\left(\hat{\psi}_{it}^{-k(i)}\left(\underline{\pi}_{t}\right)- \hat{\psi}_{i,t-1}^{-k(i)}\left(\underline{\pi}_{t-1}\right)\right)\cdot \widehat{Q}_{t}^{\pi_{(t+1):T},-k(i)}\left(H_{it},\pi_t\right)\right) \nonumber
\end{align}
and $\hat{\psi}_{it}^{-k(i)}(\underline{\pi}_{t}) \equiv \left(\prod_{s=1}^{t}\mathbf{1}\left\{ A_{is}=\pi_{s}(H_{is})\right\}\right)/(\prod_{s=1}^{t}\hat{e}_{t}^{-k(i)}\left(H_{is},A_{is}\right))$, where $\widehat{W}^{AIPW}(\pi)$ is an AIPW estimator of $W(\pi)$. This approach was originally proposed by \citet{Zhang_et_2013} without cross-fitting and under binary treatment at each stage, but its statistical properties were not established.\footnote{\cite{Jiang_Li_2016}, \cite{Thomas_et_al_2016}, and \cite{Kallus_Uehara_2020} propose AIPW estimators of welfare functions for evaluating fixed DTRs, but their focus is not on optimizing DTRs.}
One advantage of this method is that it does not require the correct specification of $\Pi$ (Assumption \ref{asm:first-best}) for consistent estimation of the optimal DTR.

However, this approach faces two computational challenges. First, since the nuisance components $\{Q_{t}^{\pi_{(t+1):T}}\}_{t=1,\ldots,T}$ depend on each specific DTR $\pi$, implementing the method requires estimating $\{Q_{t}^{\pi_{(t+1):T}}\}_{t=1,\ldots,T}$ for every candidate DTR $\pi$.\footnote{A heuristic alternative to solve this difficulty, proposed by \cite{Zhang_et_2013}, is to first estimate the optimal policy by Q-learning, denoted by $\hat{\pi}_{t}^{Q}$, and then use the optimal Q-function estimate $\widehat{Q}_{t}^{\hat{\pi}_{(t+1):T}^{Q},-k(i)}\left(H_{it},a_{t}\right)$ instead of $\widehat{Q}_{t}^{\pi_{(t+1):T},-k(i)}\left(H_{it},a_t\right)$ when constructing and maximizing $\widehat{W}^{AIPW}(\pi)$. This approach avoids the need to estimate $Q_{t}^{\pi_{(t+1):T}}$ for each $\pi_{(t+1):T}$.
} \citet{Nie_et_al_2021} highlight this computational burden.
Second, maximizing $\widehat{W}^{AIPW}(\pi)$ simultaneously across all stages is computationally demanding, especially when $T$ is not very small, as the problem is non-convex. These computational issues make the approach intractable unless the class of DTRs $\Pi$ is small (e.g., consisting of a limited number of candidate policies).

In contrast, our proposed approach offers notable computational advantages: (i) the nuisance components $\{Q_{t}^{\hat{\pi}_{(t+1):T}}\}_{t=1,\ldots,T}$ depend only on the previously estimated policies $\hat{\pi}_{(t+1)}$ from earlier steps in the backward optimization, and (ii) the backward optimization is computationally more efficient than joint optimization across all stages. While our approach requires estimating the Q-function $Q_{t}^{\hat{\pi}_{(t+1):T}}$ at each step of the sequential procedure, the associated computational cost is relatively modest.

Although the simultaneous maximization approach (\ref{eq:simultaneous_maximization}) is not our primary proposal, we show its statistical properties as follows.\footnote{\citeauthor{Sakaguchi_2021} (\citeyear{Sakaguchi_2021}, Theorem E.1) shows statistical properties of this approach in the case of binary treatment.}


\medskip

\begin{theorem}\label{thm:theorem_simlutaneous}
Under Assumptions \ref{asm:sequential independence}--\ref{asm:overlap}, \ref{asm:rate_of_convergence_backward}, and \ref{asm:bounded entropy},
\begin{align}
    R(\hat{\pi}^{AIPW}) = O_{p}\left(\kappa(\Pi) \cdot n^{-1/2}\right) + O_{p} (n^{-\min\{1/2,\tau/2\}}). \label{eq:theorem_result_AIPW}
\end{align}
\end{theorem}

\begin{proof}
See Appendix \ref{app:statistical_property_AIPW}.
\end{proof}

\medskip

This theorem shows that $\hat{\pi}^{AIPW}$ attains the same convergence rate for regret as our proposed approach, $\hat{\pi}$, under the same conditions regarding the MSE convergence rate of the nuisance component estimators and the complexity of the DTR class. Since the simultaneous optimization approach does not require the correct specification of the DTR class, Theorem \ref{thm:theorem_simlutaneous} does not rely on Assumption \ref{asm:first-best}.

%%%%%%%%%%%%%%%%%%%%%%%%%%%%%%%%%%%%%%%%%%%%%%%%%%%
%%%%%%%%%%%%%%%%%%%%%%%%%%%%%%%%%%%%%%%%%%%%%%%%%%%

\section{Simulation Study}\label{sec:simulation}


We conduct a simulation study to examine the finite sample performance of the approach presented in Section \ref{sec:doubly_robust_policy_learning}.
We consider two data generating processes (DGPs), labeled DGP1 and DGP2, each of which consists of two stages of binary treatment assignment $(A_{1},A_{2}) \in \{0,1\}^2$, associated second-stage potential outcomes $\left\{Y_2\left(a_{1},a_{2}\right)\right\}_{\left\{ a_{1},a_{2}\right\} \in\left\{ 0,1\right\} ^{2}}$,
$20$ state variables $(S_{1}^{(1)},\ldots,S_{1}^{(20)})$ observed at the first stage, and one state variable $S_2$ observed at the second stage. Each DGP is structured as follows:
\begin{align*}
&(S_{1}^{(1)},\ldots,S_{1}^{(20)})^{\prime} \thicksim N(\boldsymbol{0},I_{20});\\
&S_{2}\left(a_{1}\right)=  \sign\left(S_{1}^{(1)}\right)\cdot a_{1}+ S_{1}^{(2)} + \left(S_{1}^{(3)}\right)^2 + S_{1}^{(4)} + \varepsilon_{1} \mbox{\ with\ }\varepsilon_{1} \sim N(0,1);\\
&Y_{2}\left(a_{1},a_{2}\right) =  \phi(a_1,S_2(a_1))\cdot a_{2} 
    + 0.5 \cdot S_{2}(a_1) + S_{1}^{(4)} - \left(S_{1}^{(5)}\right)^2 + S_{1}^{(6)} + \varepsilon_{2} \mbox{\ with\ } \varepsilon_{2} \sim N(0,1);\\
&A_{1} \thicksim Ber\left(1/(1+e^{0.5 S_{1}^{(2)}-0.5S_{1}^{(3)}-S_{1}^{(5)}})\right),\ A_{2} \thicksim Ber\left(1/(1+e^{0.5S_{1}^{(5)}+0.5 S_{2}-0.2A_{1}})\right).
\end{align*}
In DGPs 1 and 2, we specify $\phi(a_1,S_2(a_1))$ as $\phi(a_1,S_2(a_1)) = \sign(S_2(a_1)\cdot (a_1 - 1/2))$ and $\phi(a_1,S_2(a_1)) = S_2(a_1) + (a_1 - 1/2)$, respectively. In each DGP, the first-stage treatment $a_1$ influences the outcome $Y_2$ through both direct and indirect channels: (i) a direct effect on $Y_2$ via treatment effect heterogeneity $Y_2(a_1,1)-Y_2(a_1,0)$; (ii) an indirect effect on $Y_2$ via the second-stage state $S_2(a_1)$. 

We compare the performance of the proposed approach (labeled ``DR'') with those of three existing methods:  Q-learning without and with policy search (labeled ``Q-learn'' and ``Q-search,'' respectively) and the IPW classification-based approach with backward optimization (labeled ``IPW'').\footnote{Following \cite{Murphy_2005}, we separately consider Q-learning with and without policy search. In Q-learning with policy search, the optimal policy for each stage $t$ is chosen from a pre-specified policy class $\Pi_t$, specifically estimated as $\hat{\pi}_{t} = \argmax_{\pi_{t}\in \Pi_{t}}\sum_{i=1}^{n}\widehat{Q}_{t}^{\hat{\pi}_{(t+1):T}}(H_{it},\pi_{t})$. 
In contrast, Q-learning without policy search optimizes the policy for each stage $t$ over all measurable policies, such as $\hat{\pi}_{t}(h_t) = \argmax_{a_{t}\in \MA_{t}}\sum_{i=1}^{n}\widehat{Q}_{t}^{\hat{\pi}_{(t+1):T}}(h_{t},a_{t})$ for any $h_t$. This approach consistently estimates the first-best DTR unless the Q-functions are misspecified.}
For each method, we use generalized random forests \citep{Athey_et_al_2019} to estimate nuisance components. We set $K=5$ for cross-fitting in the proposed approach.



For DR, Q-search, and IPW, we use a class of DTRs $\Pi=\Pi_{1} \times \Pi_{2}$ with $\Pi_1$ being the class of depth-1 decision trees of $H_1$, and $\Pi_2$ being the class of depth-2 decision trees of $H_2$. In DGP1, $\Pi_2$ is correctly specified in the sense of Assumption \ref{asm:first-best}, whereas in DGP2, $\Pi_2$ is misspecified, potentially leading to a loss of optimality in backward optimization. Note that Q-learn consistently estimates the optimal DTRs in both DGPs. We solve the optimization problems involving decision trees using the exact learning algorithm proposed by \cite{zhou2022offline}.

%\smallskip

%\begin{center}
%[Table~\ref{table:simulatin_results_dgp1} about here]
%\end{center}

%\smallskip

%\begin{center}
%[Table~\ref{table:simulatin_results_dgp2} about here]
%\end{center}


%\smallskip

Tables \ref{table:simulatin_results_dgp1} and \ref{table:simulatin_results_dgp2} present the results of 500 simulations with sample sizes of $n = 250$, $500$, $1000$, $2000$, and $4000$ for DGPs 1 and 2, respectively. In each simulation, welfare is calculated using a test sample of 50,000 observations randomly drawn from the same DGP. The results show that DR consistently outperforms the other methods in terms of mean welfare across all sample sizes for both DGPs. For example, in DGP1 with a small sample size ($n=500$), DR achieves over 40\% higher welfare than any other method. Notably, DR surpasses Q-learn even in DGP2, where the DTR class is misspecified, although Q-learn consistently estimates the optimal DTR.



Appendix \ref{app:additional_simulation_results} presents additional simulation results that examine the effects of misspecification of either the Q-functions or  propensity scores. The results demonstrate the doubly robust property of the proposed method.
%\medskip




%%%%%%%%%%%%%%%%%%%%%%%%%%%%%%%%%%%%%%%%%%%%%%%%%%%

\section{Empirical Application}\label{sec:empirical_illustration}

We apply the proposed approach to data from Project STAR 
\citep[e.g.,][]{krueger_1999, Ding_Lehrer_2010}, 
where we study the optimal allocation of students to regular-size classes (22 to 25 students per teacher) with a full-time teacher aide and small-size classes (13 to 17 students per teacher) without one in their early education (kindergarten and grade 1).\footnote{\cite{krueger_1999} reports that the presence of a teacher aide did not have a significant impact on student test scores. However, whether teaching aides have effects on academic attainment has not been examined by accounting for multiple stages of treatment and treatment effect heterogeneity.} We use a dataset of 1,877 students who were assigned to either regular-size classes with a full-time teacher aide or small-size classes without one in kindergarten.\footnote{We exclude regular-size classes without a teacher aide, as they are unlikely to be preferable to either regular-size classes with an aide or small-size classes without one for any student.} Among these students, 702 were randomly assigned to regular-size classes with a teacher aide, while the remaining students were assigned to small-size classes without a teacher aide in kindergarten (labeled ``grade K''). Upon their progression to grade 1, students were expected to remain in the same class type. However, about 10\% of students switched class types on their own (see, e.g., \cite{Ding_Lehrer_2010} for a detailed discussion). We leverage this variation to estimate the optimal DTR and consider this empirical task in the observational data setting.

We investigate the optimal allocation of students to the two class types in grades K and 1, based on their socioeconomic backgrounds, educational environment, and intermediate academic achievement. Each student’s academic achievement is measured by their percentile rank on combined reading and mathematics test scores taken at the end of grade 1. The welfare objective for the administrator is assumed to be maximizing the population average of this academic achievement measure. 

We define the first and second stages ($t=1$ and $t=2$) as grades K and 1, respectively. 
We define the action set $\MA_{t} = \{\text{aide},\text{small}\}$,
where the treatment variable $A_{t}$ is labeled ``aide'' if a student is in a regular-size class with a teacher aide at stage $t$, and ``small'' if the student is in a small-size class. The outcome variable $Y_2$ denotes the percentile rank of the combined reading and mathematics scores at the end of grade 1.
We do not incorporate the first-stage outcome into the objective function $W(\pi)$.

We use seven variables in $H_1 (=S_{1})$: student gender, student ethnicity (White/Asian or other), eligibility for free or reduced-price school lunch, school location type (rural or non-rural), teacher’s degree (bachelor’s or higher), years of teaching experience, and teacher ethnicity (White or other). For the second-stage state variables $S_{2}$, we include three variables: reading, math, and total test scores at the end of kindergarten. Recall that $H_2 = (A_1, S_1, S_2)$. We assume that the academic and socio-economic information contained in $H_2$ satisfies the conditional ignorability assumption (Assumption \ref{asm:first-best}) for grade 1, as these factors are strongly associated with the self-selection of class type.  Since there is no self-selection in kindergarten, the ignorability condition for the first stage is assured.


The class of DTRs $ \Pi = \Pi_1 \times \Pi_2$ is defined as follows. For the policy class $\Pi_1$ associated with class allocation in grade K, we employ a class of depth-1 trees that may take splitting variables from teacher degree, teacher experience, and school location type. For the policy class $\Pi_2$ associated with class allocation in grade 1, we use a class of depth-2 trees that may take splitting variables from reading, math, and total test scores at the end of kindergarten, as well as the kindergarten class type. Note that we exclude student gender, student ethnicity, and teacher ethnicity as splitting variables, as using them for treatment choice would be discriminatory. 


In applying the proposed approach, we employ 5-fold cross-fitting and use generalized random forests \citep{Athey_et_al_2019} to estimate the nuisance components. The decision trees are optimized using the exact learning algorithm of \cite{zhou2022offline}.

%\smallskip

%\begin{center}
%[Figure~\ref{fig:estimated_decision_trees} about here]
%\end{center}

%\smallskip

Figure \ref{fig:estimated_decision_trees} shows the DTR estimated using the proposed approach. The policy for grade K class allocation uses teacher experience to determine the class type, indicating that teachers with 19 years of experience or less should be assigned to small-size classes. The policy for grade 1 class allocation uses total and reading test scores at the end of kindergarten to determine each student's class type. For instance, under the estimated policy, students with a total test score below 914 are assigned to the small-size class in grade 1.

%\smallskip

%\begin{center}
%[Table~\ref{table:empirical_results} about here]
%\end{center}

%\smallskip

Table \ref{table:empirical_results} presents estimates of the welfare contrasts for the optimal DTR, $\pi^{\ast,opt}$, relative to uniformly assigning all students to either aides or small classes in both grades (i.e., $W(\pi^{\ast,opt}) - \E[Y_2(\text{aide},\text{aide})]$ and $W(\pi^{\ast,opt}) - \E[Y_2(\text{small},\text{small})]$). The table also reports the proportion of students allocated to each arm $(a_1,a_2) \in \{\text{aide},\text{small}\}^2$. These estimates were obtained using 5-fold cross-validation. The results show that class-type allocations under the optimal DTR improve academic achievement by 8.16\% and 1.27\%, respectively, compared to uniform allocations to aides or small-size classes. Additionally, under the DTR allocation, around 23\% of students are placed in regular-size classes with a teacher aide in either grade K or 1. Given the higher cost of small-size classes relative to regular-size classes with a teacher aide on a per-student basis \citep{Word_et_al_1990}\footnote{According to \cite{Word_et_al_1990}, adding a full-time aide in Grades K-3 across Tennessee cost approximately 75 million dollars annually, while reducing class sizes by one-third cost around 196 to 205 million dollars per year.}, this finding suggests that allocation based on the DTR can reduce the costs associated with class-size reduction while simultaneously improving students' academic achievement.


%%%%%%%%%%%%%%%%%%%%%%%%%%%%%%%%%%%%%%%%%%%%%%%%%%%
%\begin{comment}
\section{Conclusion}\label{sec:conclusion}
We studied the statistical learning of the optimal DTR using observational data and developed a novel doubly robust approach for learning it under the assumption of sequential ignorability. Based on backward induction, the approach learns the optimal DTR sequentially, ensuring computational tractability. Our main result shows that the resulting DTR achieves the optimal convergence rate of $n^{-1/2}$ for welfare regret under mild conditions on the MSE convergence rate for estimators of the propensity scores and Q-functions. The simulation study confirms that the proposed approach outperforms other methods in finite sample settings. Applying the proposed approach to Project STAR data, we estimate the optimal DTR for the sequential allocation of students to regular-size classes with a teacher aide and small-size classes without one in early education.
%\end{comment}



%%%%%%%%%%%%%%%%%%%%%%%%%%%%%%%%%%%%%%%%%%%%%%%%%%%

\newpage
\appendix
\section*{Tables}






\begin{table}[h]
\centering \caption{Monte Carlo simulation results for DGP1}
\label{table:simulatin_results_dgp1}
\medskip
\scalebox{1}{
\begin{tabular}{cccccc}
\hline 
& \multicolumn{5}{c}{Sample size ($n$)} \\
\cmidrule(l{3pt}r{3pt}){2-6} 
Method  & 250 & 500 & 1000 & 2000 & 4000 \\ 
\hline
\multirow{2}{*}{Q-learn} & 0.20 & 0.31 & 0.43 & 0.57 & 0.68 \\ 
   & (0.12) & (0.11) & (0.10) & (0.07) & (0.03) \\ 
\multirow{2}{*}{Q-search} & 0.17 & 0.29 & 0.40 & 0.53 & 0.64 \\ 
   & (0.19) & (0.21) & (0.21) & (0.16) & (0.10) \\ 
\multirow{2}{*}{IPW} & 0.24 & 0.29 & 0.37 & 0.45 & 0.55 \\ 
   & (0.13) & (0.14) & (0.16) & (0.14) & (0.13) \\ 
\multirow{2}{*}{DR} & 0.29 & 0.44 & 0.61 & 0.70 & 0.72 \\ 
   & (0.17) & (0.17) & (0.11) & (0.04) & (0.02) \\ 
\hline 
\end{tabular}


}
\begin{tablenotes} \footnotesize
\item Notes: Each cell shows the mean welfare, with the standard deviation in parentheses, for each method and sample size. These values are calculated based on 500 simulations using a test sample of 50,000 observations randomly drawn from DGP1.
\end{tablenotes} 
\end{table}


\begin{table}[h]
\centering \caption{Monte Carlo simulation results for DGP2}
\label{table:simulatin_results_dgp2}
\medskip
\scalebox{1}{
\begin{tabular}{cccccc}
\hline 
& \multicolumn{5}{c}{Sample size ($n$)} \\
\cmidrule(l{3pt}r{3pt}){2-6} 
Method  & 250 & 500 & 1000 & 2000 & 4000 \\ 
\hline
\multirow{2}{*}{Q-learn} & 1.31 & 1.69 & 1.78 & 1.84 & 1.89 \\ 
   & (0.36) & (0.11) & (0.08) & (0.05) & (0.03) \\ 
\multirow{2}{*}{Q-search} & 1.26 & 1.58 & 1.54 & 1.52 & 1.54 \\ 
   & (0.48) & (0.18) & (0.09) & (0.04) & (0.08) \\ 
\multirow{2}{*}{IPW} & 1.30 & 1.42 & 1.56 & 1.68 & 1.81 \\ 
   & (0.21) & (0.17) & (0.19) & (0.21) & (0.19) \\ 
\multirow{2}{*}{DR} & 1.51 & 1.77 & 1.95 & 2.01 & 2.04 \\ 
   & (0.26) & (0.24) & (0.14) & (0.07) & (0.03) \\ 
\hline 
\end{tabular}


}
\begin{tablenotes} \footnotesize
\item Notes: Each cell shows the mean welfare, with the standard deviation in parentheses, for each method and sample size. These values are calculated based on 500 simulations using a test sample of 50,000 observations randomly drawn from DGP2.
\end{tablenotes} 
\end{table}





\begin{comment}

\begin{table}[h]
\centering \caption{Monte Carlo simulation results for DGP1}
\label{table:simulatin_results_dgp1}
\medskip
\scalebox{1}{
\begin{tabular}{cccccc}
\hline 
& \multicolumn{5}{c}{Sample size ($n$)} \\
\cmidrule(l{3pt}r{3pt}){2-6} 
Method  & 250 & 500 & 1000 & 2000 & 4000 \\ 
  \hline
\multirow{2}{*}{Q-learn} & 0.20 & 0.31 & 0.44 & 0.58 & 0.69 \\ 
   & (0.11) & (0.11) & (0.10) & (0.07) & (0.03) \\ 
  \multirow{2}{*}{Q-search} & 0.17 & 0.29 & 0.41 & 0.54 & 0.65 \\ 
   & (0.19) & (0.22) & (0.21) & (0.16) & (0.10) \\ 
  \multirow{2}{*}{IPW} & 0.23 & 0.30 & 0.37 & 0.46 & 0.55 \\ 
   & (0.13) & (0.15) & (0.16) & (0.15) & (0.12) \\ 
  \multirow{2}{*}{DR} & 0.28 & 0.44 & 0.62 & 0.70 & 0.73 \\ 
   & (0.18) & (0.17) & (0.10) & (0.04) & (0.02) \\ 
\hline 
\end{tabular}

}
\begin{tablenotes} \footnotesize
\item Notes: Each cell shows the mean welfare, with the standard deviation in parentheses, for each method and sample size. These values are calculated based on 500 simulations using a test sample of 50,000 observations randomly drawn from DGP1.
\end{tablenotes} 
\end{table}


\begin{table}[h]
\centering \caption{Monte Carlo simulation results for DGP2}
\label{table:simulatin_results_dgp2}
\medskip
\scalebox{1}{
\begin{tabular}{cccccc}
\hline 
& \multicolumn{5}{c}{Sample size ($n$)} \\
\cmidrule(l{3pt}r{3pt}){2-6} 
Method  & 250 & 500 & 1000 & 2000 & 4000 \\ 
  \hline
\multirow{2}{*}{Q-learn} & 1.29 & 1.69 & 1.77 & 1.82 & 1.88 \\ 
   & (0.36) & (0.13) & (0.07) & (0.05) & (0.03) \\ 
  \multirow{2}{*}{Q-search} & 1.27 & 1.60 & 1.53 & 1.52 & 1.53 \\ 
   & (0.45) & (0.21) & (0.10) & (0.06) & (0.09) \\ 
  \multirow{2}{*}{IPW} & 1.32 & 1.41 & 1.52 & 1.66 & 1.81 \\ 
   & (0.20) & (0.18) & (0.19) & (0.21) & (0.19) \\ 
  \multirow{2}{*}{DR} & 1.50 & 1.76 & 1.93 & 1.99 & 2.03 \\ 
   & (0.26) & (0.24) & (0.14) & (0.07) & (0.03) \\ 
\hline 
\end{tabular}

}
\begin{tablenotes} \footnotesize
\item Notes: Each cell shows the mean welfare, with the standard deviation in parentheses, for each method and sample size. These values are calculated based on 500 simulations using a test sample of 50,000 observations randomly drawn from DGP2.
\end{tablenotes} 
\end{table}



\end{comment}




%--------------------%
\begin{table}[h!]
\begin{centering}
\caption{Empirical results for optimal class-type allocation}
\medskip
\label{table:empirical_results}
\scalebox{0.95}{
\begin{tabular}{cccccc}
\hline 
 \multicolumn{2}{c}{Welfare contrast}  & \multicolumn{4}{c}{Share of students in each allocation arm $(a_1,a_2)$}   \\  
 \cmidrule(r){1-2} \cmidrule(r){3-6}
 (aide, aide) & (small, small) 
 & (aide, aide)  & (small, aide)  & (aide, small) & (small, small) \\
\cmidrule(r){1-2} \cmidrule(r){3-6} 
8.16\% & 1.27\% & 1.0\%  & 17.2\%  & 5.1\% & 76.7\%\\
\hline 
\end{tabular}
}
\par\end{centering}
\centering{}
\begin{tablenotes} 
\footnotesize 
\item Notes: The first and second columns present the estimates of the welfare contrasts, defined as $W(\pi^{\ast,opt}) - \E[Y_2(\text{aide},\text{aide})]$ and $W(\pi^{\ast,opt}) - \E[Y_2(\text{small},\text{small})]$, respectively. The third through sixth columns show the estimated shares of students assigned to the four allocation arms, $(a_1, a_2) \in \{\text{aide}, \text{small}\}^2$, by the optimal DTR.
\end{tablenotes} 
\end{table}
%\bigskip
%--------------------%

%--------------------%
\begin{comment}
\begin{table}[h]
\begin{centering}
\caption{Empirical results for optimal class-type allocation}
\medskip
\label{table:empirical_results}
\begin{tabular}{ccccc}
\hline \hline
\multirow{2}{*}{Welfare contrast}  & \multicolumn{4}{c}{Share of students in each allocation arm $(a_1,a_2)$}   \\ \cline{2-5}
 & (aide, aide)  & (small, aide)  & (aide, small) & (small, small) \\
\hline 
1.27\% & 1.0\%  & 17.2\%  & 5.1\% & 76.7\%\\
\hline \hline
\end{tabular}
\par\end{centering}
\centering{}
\begin{tablenotes} 
\footnotesize 
\item Notes: The first column shows the estimate of the welfare contrast, defined as $W(\pi^{\ast,opt}) - \E[Y_2(\text{small},\text{small})]$. Each of the second to fifth columns shows the estimated share of students assigned to each of the four allocation arms, $(a_1, a_2) \in \{\text{aide}, \text{small}\}^2$, by the optimal DTR.
\end{tablenotes} 
\end{table}
%--------------------%



%--------------------%
\begin{table}[h]
\begin{centering}
\caption{Empirical results for optimal class-type allocation}
\medskip
\label{table:empirical_results}
\scalebox{0.9}{
\begin{tabular}{cccccc}
\hline \hline
&\multirow{2}{*}{Welfare contrast}  & \multicolumn{4}{c}{Share of students in each allocation arm $(a_1,a_2)$}   \\ \cline{3-6}
& & (aide, aide)  & (small, aide)  & (aide, small) & (small, small) \\
\hline 
vs. (aide,aide) &7.80\% & \multirow{2}{*}{1.0\%}  & \multirow{2}{*}{17.2\%}  & \multirow{2}{*}{5.1\%} & \multirow{2}{*}{76.7\%}\\
vs. (small,small) & 0.88\%  &   &   &  & \\
\hline \hline
\end{tabular}
}
\par\end{centering}
\centering{}
\begin{tablenotes} 
\footnotesize 
\item Notes: The first column shows the estimate of the welfare contrast, defined as $W(\pi^{\ast,opt}) - E[Y_2(\text{small},\text{small})]$. Each of the second to fifth columns shows the estimated share of students assigned to each of the four allocation arms, $(a_1, a_2) \in \{\text{aide}, \text{small}\}^2$, by the optimal DTR.
\end{tablenotes} 
\end{table}
\end{comment}
%--------------------%

\bigskip
\bigskip
\bigskip
\newpage
%%%%%%%%%%%%%%%%%%%%%%%%%%%%%%%%%%%%%%%%%%%%%%%%%%%
\newpage
\appendix
\section*{Figure}


\medskip

\begin{figure}[h]
\caption{Estimated DTR for class-type allocation in grades K and 1} \label{fig:estimated_decision_trees}
\vspace{.5cm}
\begin{minipage}[t]{0.44\hsize}
\begin{center}

{\small  (a) Policy for grade K} \\
\vspace{0.5cm}
\input{figures/depth1_tree_1st.tikz} 

\end{center}
\end{minipage}
\begin{minipage}[t]{0.44\hsize}
\begin{center}
{\small (b) Policy for grade 1} \\
\vspace{.5cm}
\input{figures/depth2_tree_1st.tikz}
\end{center}
\end{minipage}

\vspace{.3cm}
\begin{tablenotes}\footnotesize
\item[] Notes: This figure illustrates the estimated DTR from Section \ref{sec:empirical_illustration}. Panels (a) and (b) display the estimated policy trees for the class-type allocation in grades K and 1, respectively.
\end{tablenotes}
\end{figure}
\medskip



%%%%%%%%%%%%%%%%%%%%%%%%%%%%%%%%%%%%%%%%%%%%%%%%%%%%%%%%%%%%%%%%%%%%%%%



